\section{Improvement Attempts}\label{sec:improve}

We tested three ways of improving faithfulness and asked whether a gain at one level transfers to the others. Every
improved SAE was trained by us.

\subsection{Truth-focused fine-tuning: a frozen-reader gain only}\label{sec:finetune}

Method A fine-tunes an existing SAE for 500 steps on reconstruction of truth statements from city group A, with OWT
replay, against an OWT-only control trained for the same number of steps. A first truth-aware objective was found to be
circular, because its training teacher was also the evaluation probe; every result here uses probes fitted on original
activations of group B, which no objective sees (Appendix~\ref{app:details}). In the pre-registered experiment at layers
3, 8 and 9 (81 fine-tunes covering Method A, the OWT-only control and an output-KL variant), Method A passed every
guard, cut reconstruction error on truth statements by 61--85\,\% and lowered the pre-registered truth gap, measured by a frozen nonlinear probe, on all three seeds in all nine
layer$\times$architecture cells. The criterion of consistent improvement against both the base SAE and the control was
met in four of nine cells (outcome class ``preliminary improvement'').

This gain is specific to that frozen reader. A frozen logistic probe fitted on the same activations, and re-trained
readers on truth and on persona prompts, show no consistent change in any cell (Appendix Table~\ref{tab:impl_readers}).
Method A therefore mostly moves reconstructions back towards the region where the original probe was fitted; it does
not measurably restore information that a re-trained reader can recover. It also left causal preservation at the
layer-8 site unchanged: after every fine-tune the reconstruction still carried 0.61--0.69 of the effect (Appendix
Table~\ref{tab:e3}). Adding an output-KL term improved behaviour on every seed in every cell, but for ReLU and JumpReLU
at layers 8 and 9 it raised L0 by 5.3--8.1\,\%, beyond the pre-registered $\pm5\,\%$ guard, so those four cells are
``not creditable''. It never improved the truth gap, and its only causal gains (about $+0.02$ of the effect)
occurred in cells that violated the sparsity guard.

\subsection{More data and wider SAEs: confirmed reconstruction and behaviour gains}\label{sec:r2}

Our base SAEs were trained on 101{,}600 tokens, and their validation EV was still rising when training stopped. We
therefore retrained SAEs at layers 1, 8 and 11 with width 12{,}288 and 100M tokens each, with comparisons and their rule
fixed before training, and then evaluated every improvement-phase comparison once on the OWT and truth test splits
under a frozen protocol (81 SAEs). The reconstruction and behaviour gains are \evlabel{confirmed test} findings
(Table~\ref{tab:m5}): across layers and comparisons NMSE fell by 31--74\,\% and spliced $\Delta$CE by 35--79\,\%. For
TopK at layers 8 and 11, test NMSE fell from 0.160 and 0.185 to 0.094 and 0.111, and $\Delta$CE from 0.133 and 0.204
to 0.062 and 0.096 nats. On the same held-out test, the remaining TopK splice damage is still 1.20 and 1.48 times
that of norm-matched random error at layers 8 and 11, so the remaining error is still more harmful than random error
of the same size \citep{engels2024darkmatter}.

Truth information did not follow. No comparison improved the re-trained truth gap or the SAE-minus-PCA gap on test,
and the width ablation made both slightly worse at layer 8 ($|d|=0.012$--0.015 AUROC, about one quarter of the
minimum detectable effect); these are small regressions, not robust evidence of a meaningful degradation. This is an
absence of evidence within the design's sensitivity: differences between SAE families smaller than about 0.05--0.09
AUROC could not have been detected, and the observed differences are at most 0.025. The re-trained truth gap also does
not discriminate trained SAEs from the frozen-random-decoder baseline at layer 8, where on validation data the baseline
had the smaller gap (0.052 versus 0.081 for the matched trained SAE, seed 42). On validation data, the re-trained
persona gap fell by 27--31\,\% for TopK at layers 8 and 11, but the persona gap relative to error-matched PCA did not
improve, and causal preservation rose from 0.62 to 0.73 at layer 8 and from 0.62 to 0.69 at layer 11. A single
pre-registered evaluation on the held-out RepE-LM test split then tested these representation and causal gains. The
per-pair persona direction cosine improved in all nine cells (three comparisons at layers 1, 8 and 11;
\evlabel{confirmed test}). The causal-preservation and persona re-fit hypotheses were not confirmed as a whole: only
their TopK layer-8 cells passed, with causal preservation rising from 0.602 to 0.729 (Holm $p=0.021$) and the persona
re-fit gap changing by $-0.048$ (Holm $p=0.046$). Both cells are fragile: once statement-sampling uncertainty is
included in a sensitivity analysis, their Holm $p$ is 0.156. The layer-11 causal gain remains \evlabel{validation}
only.

\subsection{Behaviour-targeted training: a trade-off}\label{sec:r4}

Because the earlier output-KL gains were confounded by sparsity drift, we fine-tuned the scaled TopK SAEs, whose L0
is fixed at 128, with a reconstruction plus output-KL objective \citep{karvonen2025kl,braun2024e2e} for 5M tokens at
layers 8 and 11, against a reconstruction-only control on identical tokens. On the held-out test (\evlabel{confirmed
test}), spliced $\Delta$CE fell from 0.062 and 0.096 to 0.051 and 0.081 nats, an improvement of 16--17\,\%, while NMSE
rose by about 2\,\%, a regression within the pre-registered explained-variance guard. On validation data no information
gap changed, the per-pair persona direction cosine fell by 0.012--0.013, and causal preservation at layer 11 decreased
from 0.694 to 0.671 on every seed (outcome class ``regression''); at layer 8 it did not change (Appendix Table~\ref{tab:r4}). A
behavioural objective thus improves behaviour creditably once sparsity is fixed, at a small cost in reconstruction and
a loss of causal preservation at one of two sites.

% generated by scripts/shared/build_paper_artifacts.py from results/shared/m5_confirmatory_test/comparisons.csv
\begin{table}[t]
\centering
\small
\caption{M5: the single pre-registered confirmatory TEST pass (OWT test, 100 sequences, 12,700 tokens; GoT cities+neg\_cities test). Same rules as on validation (C/A: all 9 seed pairs and Welch CI; K: all 3 paired seeds and $|\bar d|>2\,$sd). Means over 3 seeds. C1: TopK old-3072 $\rightarrow$ R2-12288; C3: old JumpReLU-3072 $\rightarrow$ R2 BatchTopK-12288; K1: R4 KL fine-tune vs MSE control. All comparisons are counted in Table~\ref{tab:m5_counts}. Source: \texttt{results/\allowbreak{}shared/\allowbreak{}m5\_confirmatory\_test/\allowbreak{}comparisons.csv}.}
\label{tab:m5}
\begin{adjustbox}{max width=\linewidth}
\begin{tabular}{lllcccl}
\toprule
ID & Layer & Metric & Before & After & Improvement & Class (TEST) \\
\midrule
C1 & L1 & OWT NMSE $\downarrow$ & 0.0787 & 0.0248 & +68.6\% & IMPROVEMENT \\
 & L1 & spliced $\Delta$CE $\downarrow$ & 0.0221 & 0.0056 & +74.9\% & IMPROVEMENT \\
 & L8 & OWT NMSE $\downarrow$ & 0.1598 & 0.0941 & +41.1\% & IMPROVEMENT \\
 & L8 & spliced $\Delta$CE $\downarrow$ & 0.1330 & 0.0620 & +53.4\% & IMPROVEMENT \\
 & L8 & GoT re-fit gap $\downarrow$ & 0.0607 & 0.0558 & +8.1\% & no clear change \\
 & L11 & OWT NMSE $\downarrow$ & 0.1848 & 0.1112 & +39.8\% & IMPROVEMENT \\
 & L11 & spliced $\Delta$CE $\downarrow$ & 0.2041 & 0.0962 & +52.9\% & IMPROVEMENT \\
 & L11 & GoT re-fit gap $\downarrow$ & 0.1009 & 0.0775 & +23.2\% & no clear change \\
\midrule
C3 & L1 & OWT NMSE $\downarrow$ & 0.0675 & 0.0178 & +73.7\% & IMPROVEMENT \\
 & L1 & spliced $\Delta$CE $\downarrow$ & 0.0088 & 0.0019 & +78.8\% & IMPROVEMENT \\
 & L8 & OWT NMSE $\downarrow$ & 0.1560 & 0.0897 & +42.5\% & IMPROVEMENT \\
 & L8 & spliced $\Delta$CE $\downarrow$ & 0.1212 & 0.0542 & +55.3\% & IMPROVEMENT \\
 & L8 & GoT re-fit gap $\downarrow$ & 0.0584 & 0.0654 & -12.0\% & no clear change \\
 & L11 & OWT NMSE $\downarrow$ & 0.1851 & 0.1073 & +42.0\% & IMPROVEMENT \\
 & L11 & spliced $\Delta$CE $\downarrow$ & 0.2155 & 0.0927 & +57.0\% & IMPROVEMENT \\
 & L11 & GoT re-fit gap $\downarrow$ & 0.1070 & 0.0871 & +18.6\% & no clear change \\
\midrule
K1 & L8 & OWT NMSE $\downarrow$ & 0.0938 & 0.0958 & -2.1\% & \textbf{REGRESSION} \\
 & L8 & spliced $\Delta$CE $\downarrow$ & 0.0623 & 0.0514 & +17.4\% & IMPROVEMENT \\
 & L8 & GoT re-fit gap $\downarrow$ & 0.0585 & 0.0484 & +17.3\% & no clear change \\
 & L11 & OWT NMSE $\downarrow$ & 0.1109 & 0.1133 & -2.2\% & \textbf{REGRESSION} \\
 & L11 & spliced $\Delta$CE $\downarrow$ & 0.0964 & 0.0807 & +16.3\% & IMPROVEMENT \\
 & L11 & GoT re-fit gap $\downarrow$ & 0.0780 & 0.0784 & -0.5\% & no clear change \\
\bottomrule
\end{tabular}
\end{adjustbox}
\end{table}

